\question[4]
Ein Schüler behauptet: ,,Ein Verschlüsselungsverfahren ist nur dann 
sicher, wenn der Algorithmus, also die genaue Funktionsweise der 
Verschlüsselung, geheim bleibt. Kennt der Angreifer das Verfahren, ist 
die Verschlüsselung wertlos.``

Erläutere das Kerckhoffs'sche Prinzip und nimm begründet Stellung zu 
der Aussage des Schülers.

\begin{solutionbox}{6cm}

Das Kerckhoffs'sche Prinzip besagt, dass die Sicherheit eines 
Verschlüsselungsverfahrens allein auf der Geheimhaltung des Schlüssels 
beruhen soll, nicht auf der Geheimhaltung des Algorithmus selbst. Der 
Algorithmus darf also öffentlich bekannt sein; ein Angreifer, der das 
Verfahren kennt, aber nicht den Schlüssel, soll den Klartext trotzdem 
nicht ermitteln können.

Die Aussage des Schülers ist daher nicht richtig. Ein Verfahren, dessen 
Sicherheit von der Geheimhaltung des Algorithmus abhängt, gilt als 
unsicher, denn sobald der Algorithmus bekannt wird -- etwa durch 
Reverse Engineering, Spionage oder Veröffentlichung --, ist das 
gesamte Verfahren gebrochen und muss durch ein neues ersetzt werden. 
Zudem lässt sich ein geheimer Algorithmus in der Praxis kaum dauerhaft 
geheim halten und nicht von unabhängigen Fachleuten auf Schwachstellen 
überprüfen. 

\end{solutionbox}


